\documentclass{article}
\usepackage{amsmath}
\usepackage{amssymb}
\usepackage{graphicx}
\usepackage{hyperref}
\usepackage{physics}
\usepackage{tikz}
\usepackage{bm}
\usepackage{booktabs}
\usepackage{siunitx}

\title{Derivation of Nontrivial LICI Geometries}
\date{\today}

\begin{document}

\maketitle

\section{Introduction}

This document presents the mathematical derivation for finding nontrivial Light-Induced Conical Intersection (LICI) geometries in a three-molecule system. The system consists of three identical diatomic molecules placed in a cavity with their axes parallel to the $x$, $y$, and $z$ axes of a Cartesian coordinate system.

\section{Model Potentials}

The system is described by two potential energy surfaces:
\begin{align}
V_x(x) &= x^2 \\
V_a(x) &= a \cdot (x - x_{\text{shift}})^2 + c_0
\end{align}

Where:
\begin{itemize}
    \item $a$ is a scaling parameter
    \item $x_{\text{shift}}$ is a shift parameter for the second potential
    \item $c_0$ is a constant energy offset
\end{itemize}

\section{Potential Difference Analysis}

The difference between the two potentials is:
\begin{align}
V_a(x) - V_x(x) &= a \cdot (x - x_{\text{shift}})^2 + c_0 - x^2 \\
&= (a - 1) \cdot x^2 - 2 \cdot a \cdot x_{\text{shift}} \cdot x + c_1
\end{align}

Where $c_1 = a \cdot x_{\text{shift}}^2 + c_0$.

This can be rewritten in the form of a perfect square plus a constant:
\begin{equation}
V_a(x) - V_x(x) = (a - 1) \cdot (x - x_{\text{prime}})^2 + c_1'
\end{equation}

Where $x_{\text{prime}}$ is the point where the quadratic term has its minimum:
\begin{equation}
x_{\text{prime}} = \frac{a}{a - 1} \cdot x_{\text{shift}}
\end{equation}

\section{Trivial and Nontrivial LICIs}

\subsection{Trivial LICIs}

Trivial LICIs occur when all three bond distances are equal:
\begin{equation}
(r_1, r_2, r_3) \text{ with } r_1 = r_2 = r_3 (= r_0)
\end{equation}

These form a seam with the property $r_1 = r_2 = r_3$.

\subsection{Nontrivial LICIs}

To find nontrivial LICIs, we search for points where two different bond distances $r_1 \neq r_2$ give equal potential differences:
\begin{equation}
(V_a - V_x)(r_1) = (V_a - V_x)(r_2)
\end{equation}

If such $r_1 \neq r_2$ exist, the following configurations provide nontrivial LICIs:
\begin{align}
\text{LICI}_1 &: (r_1, r_2, r_3) = (r_1, r_1, r_2) \\
\text{LICI}_2 &: (r_1, r_2, r_3) = (r_1, r_2, r_1) \\
\text{LICI}_3 &: (r_1, r_2, r_3) = (r_2, r_1, r_1) \\
\text{LICI}_4 &: (r_1, r_2, r_3) = (r_2, r_2, r_1) \\
\text{LICI}_5 &: (r_1, r_2, r_3) = (r_2, r_1, r_2) \\
\text{LICI}_6 &: (r_1, r_2, r_3) = (r_1, r_2, r_2)
\end{align}

\section{Geometric Constraints}

Since the trivial LICIs form a seam, we set up our closed paths in the plane perpendicular to this "trivial" seam. All points in this plane have the property:
\begin{equation}
r_1 + r_2 + r_3 = 3 \cdot r_0
\end{equation}
where $(r_0, r_0, r_0)$ is the trivial LICI in the plane.

It is easy to see that $\text{LICI}_1$, $\text{LICI}_2$, and $\text{LICI}_3$ lie in the same plane, and similarly, $\text{LICI}_4$, $\text{LICI}_5$, and $\text{LICI}_6$ also lie in the same plane.

\section{Derivation for LICI$_1$, LICI$_2$, LICI$_3$}

We now search for nontrivial LICIs of type $\text{LICI}_1$, $\text{LICI}_2$, $\text{LICI}_3$ with the additional requirement:
\begin{equation}
r_1 + r_1 + r_2 = 3 \cdot r_0
\end{equation}
for a predefined $r_0$.

From the equality condition $(V_a - V_x)(r_1) = (V_a - V_x)(r_2)$ with $r_1 \neq r_2$, we obtain:
\begin{equation}
r_1 = x_{\text{prime}} - \delta, \quad r_2 = x_{\text{prime}} + \delta
\end{equation}
where $\delta$ is an arbitrary parameter.

Applying the geometric constraint:
\begin{align}
r_1 + r_1 + r_2 &= 3 \cdot r_0 \\
3 \cdot x_{\text{prime}} - \delta &= 3 \cdot r_0 \\
\delta &= 3 \cdot (x_{\text{prime}} - r_0)
\end{align}

Substituting back to find $r_1$ and $r_2$:
\begin{align}
r_1 &= x_{\text{prime}} - \delta = 3 \cdot r_0 - 2 \cdot x_{\text{prime}} = r_0 - 2 \cdot (x_{\text{prime}} - r_0) \\
r_2 &= x_{\text{prime}} + \delta = 4 \cdot x_{\text{prime}} - 3 \cdot r_0 = r_0 + 4 \cdot (x_{\text{prime}} - r_0)
\end{align}

\section{Distance from Trivial LICI}

To determine how many LICIs are enclosed by a given circle, we need to know the distance between the nontrivial and trivial LICIs. The distance in the 3D coordinate space is:
\begin{align}
d_{\text{CI}} &= |r_2 - r_0| \cdot \frac{\sqrt{6}}{2} \\
&= 4 \cdot |x_{\text{prime}} - r_0| \cdot \frac{\sqrt{6}}{2} \\
&= 2 \cdot \sqrt{6} \cdot |x_{\text{prime}} - r_0|
\end{align}

\section{Analysis of LICI$_4$, LICI$_5$, LICI$_6$}

Similar expressions could be derived for $\text{LICI}_4$, $\text{LICI}_5$, and $\text{LICI}_6$. In this case, we would obtain:
\begin{equation}
\delta' = -3 \cdot (x_{\text{prime}} - r_0) = -\delta
\end{equation}

Taking this value and substituting into the expressions for $r_1$ and $r_2$:
\begin{align}
r_1' &= x_{\text{prime}} - \delta' = x_{\text{prime}} + \delta = r_2 \\
r_2' &= x_{\text{prime}} + \delta' = x_{\text{prime}} - \delta = r_1
\end{align}

Combining this with the definitions of the nontrivial $\text{LICI}_i$ configurations, we realize that these LICIs are actually the same as $\text{LICI}_1$, $\text{LICI}_2$, and $\text{LICI}_3$. More specifically:
\begin{align}
\text{LICI}_4 &= \text{LICI}_1 \\
\text{LICI}_5 &= \text{LICI}_2 \\
\text{LICI}_6 &= \text{LICI}_3
\end{align}

\section{Conclusion}

The derivation shows that for a given plane defined by $r_1 + r_2 + r_3 = 3 \cdot r_0$, there are exactly three distinct nontrivial LICI configurations. These are symmetrically arranged around the trivial LICI at $(r_0, r_0, r_0)$, with the distance from the trivial LICI given by:
\begin{equation}
d_{\text{CI}} = 2 \cdot \sqrt{6} \cdot |x_{\text{prime}} - r_0|
\end{equation}

This geometric understanding is crucial for determining the topological properties of the system and for setting up appropriate closed paths for geometric phase calculations.

\end{document}
